\subsection*{Question 1}
	Factorisons :\quad $\begin{aligned}[t]
		9x^2 - \dfrac{1}{9} &= (3x)^2 - \left(\dfrac{1}{3}\right)^2 \quad \text{on reconnaît alors une expression de la forme } a^2 - b^2\\
		&=\left(3x - \dfrac{1}{3}\right)\left(3x + \dfrac{1}{3}\right) \quad \text{qui se factorise en } (a - b)(a + b)
	\end{aligned}$

	\textbf{Bonne réponse : B.}

\subsection*{Question 2}
	Effectuons les calculs : \quad $E = \dfrac{x - y}{zt} = \dfrac{3 - (-2)}{(-3)\times(-4)} = \dfrac{3 + 2}{12} = \dfrac{5}{12}$.

	\textbf{Bonne réponse : C.}

\subsection*{Question 3}
On n'a pas besoin de lectures graphiques très précises pour déterminer le signe de $A$.

On voit \quad $f(-4) \approx 1,5$,\quad donc \quad $f(-4)>0$;

et \quad $f(-1) = -2$,\quad donc \quad $f(-1) < 0$.

$A$ est donc le quotient de deux nombres de signes contraires, le numérateur étant non nul (et le dénominateur aussi, bien sûr), donc on a :\quad $A < 0$.

\textbf{Bonne réponse : B.}

\subsection*{Question 4}
	Résolvons $(I)$ :\quad $\begin{aligned}[t] -2x + 2 \geqslant 0
		&\iff -2x \geqslant -2\\
		&\iff x \leqslant \dfrac{-2}{-2}\quad \text{car on divise par } -2 < 0\\
		&\iff x \leqslant 1
	\end{aligned}$

	L'ensemble des solutions est donc $S = ]-\infty ; 1]$.

	\textbf{Bonne réponse : A.}

\subsection*{Question 5}
	Les propositions sont des fonctions polynômes du second degré, présentées sous une forme factorisée.

	On veut ici une fonction qui aura deux racines : $-3$ et 2, et dont le coefficient dominant doit être positif, car les images sont positives, sauf entre les deux racines.

	On peut donc éliminer les propositions \textbf{B.} (les racines de cette proposition sont $-2$ et $3$) et \textbf{D.} (les racines sont $-2$ et $-3$).

	Les propositions \textbf{A.} et \textbf{C.} présentent les bonnes racines, mais la proposition \textbf{A.} a un coefficient dominant négatif : si on développe, on trouvera :\quad $-x^2 -x + 6$.

	Finalement, la seule fonction qui admet le tableau de signe donné dans l'énoncé est $f$, dont une expression est :\quad $f(x) = (x - 2)(x + 3)$.

	\textbf{Bonne réponse : C.}

\subsection*{Question 6}
	Une baisse de 50\,\% se calcule en multipliant par :\quad $c_1 = 1 - \dfrac{50}{100} = 0,5$;

	la hausse de 40\,\% en multipliant par :\quad $c_2 = 1 + \dfrac{40}{100} = 1,4$;

	L'évolution globale a donc un coefficient multiplicateur de :\quad $c_g = c_1  \times c_2 = 0,5 \times 1,4 = 0,7$.

	Le taux d'évolution global est donc de :\quad $t_g = c_g - 1 = 0,7 - 1 = -0,3 = -\dfrac{30}{100}$.

	L'évolution globale du prix est donc une baisse de 30\,\%.

	\textbf{Bonne réponse : D.}

\subsection*{Question 7}
	Si 40\,\% sont des hommes, alors les femmes représentent le complément à 100\,\%, c'est-à-dire 60\,\% des adhérents de l'association.

	30 femmes représentent donc 60\,\%, donc une personne représente 2\,\% des adhérents de l'association.

	Il y a donc en tout 50 adhérents à cette association.

	\emph{Variante :} 30 femmes représentent 60\,\% des adhérents, donc 10 femmes représentent 20\,\% des adhérents. Ainsi $10 \times 5 = 50$ personnes représentent 100\,\% des adhérents.

	\textbf{Bonne réponse : D.}

\subsection*{Question 8}
	On applique les formules de puissances :
	\begin{itemize}
		\item $\dfrac{a^8}{a^{-5}} = a^{8 - (-5)} = a^{13} \neq a^3$;
		\item $\dfrac{a^{30}}{a^2} = a^{30 - 2} = a^{28} \neq a^{15}$;
		\item $\big(a^{10}\big)^3 = a^{10 \times 3} = a^{30} \neq a^{13}$;
		\item $\dfrac{a\times a^5}{a^2} = \dfrac{a^1\times a^5}{a^2} = a^{1 + 5 - 2} = a^4$.
	\end{itemize}

	\textbf{Bonne réponse : D.}

\needspace{10\baselineskip}
